\section{完整案例：数据集概览}
下面将同一组数据分别显示为样本数和总量占比。两列分母都为 125000；
标签格式不同，但相同数据的彩色条长度保持一致。

\begin{example}[text and listing,sidebyside=false]{示例 69：样本量与占比}
\begingroup
\gradbarssetup{
  max=125000, width=29mm,
  label width=4.8em, precision=0
}
\renewcommand{\arraystretch}{1.65}
\begin{tabular}{@{}lcc@{}}
  \toprule
\multicolumn{1}{c}{Split} & \multicolumn{1}{c}{Samples} & \multicolumn{1}{c}{Share of total}\\
  \midrule
  Training
    & \gradbar[theme=solid]{87500}
    & \gradbar[theme=teal,
        value format=percent,precision=1]{87500} \\
  Validation
    & \gradbar[theme=solid]{18750}
    & \gradbar[theme=teal,
        value format=percent,precision=1]{18750} \\
  Test
    & \gradbar[theme=solid]{18750}
    & \gradbar[theme=teal,
        value format=percent,precision=1]{18750} \\
  Unassigned
    & \gradbar[theme=solid]{0}
    & \gradbar[theme=teal,
        value format=percent,precision=1]{0} \\
  \bottomrule
\end{tabular}
\endgroup
\end{example}

\note{本例演示较大数值与真实零值。样本数中的分组空白仅用于可读性，
输入仍必须写成 \key{87500}，不能在数值参数中加入逗号或单位。}
